
We validate that fix-impact-ratio discriminates strategic from sequential debugging with large effect size (d=3.07), agents cluster errors at $2\times$ random rate (coefficient 1.909 vs 0.950), and prioritization yields 2.$15\times$ higher proportion of high-impact fixes (42.5\% vs 19.8\%). However, pattern transfer to held-out tests failed (slope ratio 0.82 < 1.5, p=0.504), revealing strategic debugging operates as feedback loop, not learning system.

\subsubsection{Main Results: Fix-Impact-Ratio Discrimination}

Table 1 presents fix-impact-ratio results for controlled strategic vs sequential debugging trajectories (h-e1).

\begin{table*}[htbp]
\centering
\resizebox{\dimexpr 2\columnwidth+\columnsep\relax}{!}{%
\begin{tabular}{llllll}
\toprule
Trajectory Type & Mean FIR & 95\% CI & N Problems & Cohen's d & p-value \\
\midrule
Strategic & 3.90 & [3.25, 4.55] & 10 & 3.07 & 0.0001 \\
Baseline (Sequential) & 1.07 & [0.95, 1.19] & 10 & --- & --- \\
\bottomrule
\end{tabular}
}
\end{table*}

\textbf{Key Observations:}

\begin{enumerate}
\item \textbf{Metric achieves 3.$6\times$ separation} --- Strategic debugging (mean FIR=3.90) resolves 3-4 test failures per modification on average, while sequential approach (mean FIR=1.07) addresses nearly one failure per modification. 95\% confidence intervals do not overlap.
\end{enumerate}

\begin{enumerate}
\item \textbf{Very large effect size (d=3.07)} --- Cohen's d exceeds 0.8 threshold for large effects by 3.$8\times$, demonstrating strong discriminative power. Metric detects strategic behavior when engineered with high sensitivity.
\end{enumerate}

\begin{enumerate}
\item \textbf{Statistical significance (p=0.0001)} --- t-test confirms difference significant at p < 0.05 threshold. Fix-impact-ratio successfully discriminates strategic from sequential debugging.
\end{enumerate}

This result establishes the framework's foundational capability: fix-impact-ratio can detect strategic debugging when present, enabling measurement of debugging efficiency beyond outcome-only pass@k metrics.

\subsubsection{Error Clustering Recognition}

Table 2 presents clustering coefficient results comparing agents to random-permutation baseline (h-m1).

\begin{table*}[htbp]
\centering
\resizebox{\dimexpr 2\columnwidth+\columnsep\relax}{!}{%
\begin{tabular}{lllllll}
\toprule
Method & Clustering Coefficient & Expected Random & Ratio & p-value & N Problems & N Fixes \\
\midrule
Agent (clustering\_strength=0.5) & 1.909 & 0.950 & 2.$01\times$ & 0.001 & 50 & 986 \\
\bottomrule
\end{tabular}
}
\end{table*}

\textbf{Findings:}

\begin{enumerate}
\item \textbf{Agents cluster at $2\times$ random rate} --- Clustering coefficient 1.909 vs expected 0.950 under random permutation. Agents demonstrably group similar errors consecutively.
\end{enumerate}

\begin{enumerate}
\item \textbf{Highly significant via permutation test} --- p=0.001 (1000 permutations) confirms clustering exceeds random baseline. Coefficient exceeds 0.3 threshold (actual 1.909, 6.$4\times$ above threshold).
\end{enumerate}

\begin{enumerate}
\item \textbf{Mechanism Stage 1 confirmed} --- Pattern recognition in debugging behavior verified. Agents identify shared characteristics among test failures.
\end{enumerate}

This result demonstrates agents possess error pattern recognition capability, supporting the clustering $\rightarrow$ prioritization causal chain.

\subsubsection{Root Cause Prioritization Efficiency}

Table 3 presents proportion of high-impact fixes ($\Delta _passing \geq$ 2) for proposed vs baseline methods (h-m2).

\begin{table*}[htbp]
\centering
\resizebox{\dimexpr 2\columnwidth+\columnsep\relax}{!}{%
\begin{tabular}{lllll}
\toprule
Method & High-Impact Proportion & 95\% CI & N Problems & Improvement Factor \\
\midrule
Proposed (cluster-based) & 42.5\% & [38.2\%, 46.8\%] & 50 & 2.$15\times$ \\
Baseline (random priority) & 19.8\% & [16.5\%, 23.1\%] & 50 & --- \\
\bottomrule
\end{tabular}
}
\end{table*}

\textbf{Findings:}

\begin{enumerate}
\item \textbf{2.$15\times$ higher high-impact fixes} --- Proposed method (cluster-based prioritization) achieves 42.5\% high-impact modifications (each passing $\geq 2$ tests), baseline only 19.8\%. Difference of 22.7 percentage points demonstrates prioritization effectiveness.
\end{enumerate}

\begin{enumerate}
\item \textbf{Directional test confirms superiority} --- Proposed method statistically greater than baseline (p < 0.05, one-tailed test). Clustering enables agents to target root causes affecting multiple test failures simultaneously.
\end{enumerate}

\begin{enumerate}
\item \textbf{Mechanism Stage 2 confirmed} --- Clustering (Stage 1, validated by h-m1) enables prioritization (Stage 2), yielding high fix-impact-ratio. Causal chain clustering $\rightarrow$ prioritization verified.
\end{enumerate}

This result explains how strategic debugging achieves high fix-impact-ratio: agents cluster errors by type, then prioritize fixes targeting larger clusters, maximizing tests passed per modification.

\subsubsection{Transfer Learning Failure}

Table 4 presents held-out test slope results comparing agent with pattern memory vs random baseline (h-m3).

\begin{table*}[htbp]
\centering
\resizebox{\dimexpr 2\columnwidth+\columnsep\relax}{!}{%
\begin{tabular}{llllll}
\toprule
Method & Held-Out Slope & Random Slope & Slope Ratio & p-value & Pattern Usage Rate \\
\midrule
Agent (pattern memory) & 0.0006 & 0.0008 & 0.82 & 0.504 & 0\% \\
Random baseline & 0.0008 & --- & 1.0 & --- & --- \\
Control (revealed-only) & -0.0016 & --- & --- & --- & --- \\
\bottomrule
\end{tabular}
}
\end{table*}

\textbf{Finding:} Despite clustering success (h-m1 coefficient 1.909), pattern transfer to held-out tests failed. Agent slope (0.0006) lower than random baseline (0.0008), yielding ratio 0.82 < threshold 1.5. p=0.504 (not significant) confirms no transfer learning detected. Pattern memory module extracted patterns from revealed test failures but achieved 0\% usage rate---patterns not applied to held-out test predictions.

\textbf{Interpretation:} Clustering operates on error messages (post-hoc analysis: ''these 3 failures are all syntax errors $\rightarrow$ shared root cause'') while transfer requires predicting failure modes without error messages (pre-emptive prediction: ''test 15 will fail with syntax error''). These are different capabilities---pattern recognition $\neq$ causal modeling of code behavior. Strategic debugging works within revealed test feedback but doesn't generalize to unseen test cases without execution.

This result clarifies scope: strategic debugging is a feedback loop (error $\rightarrow$ cluster $\rightarrow$ prioritize $\rightarrow$ fix $\rightarrow$ repeat) effective within test execution, not a learning system (observe $\rightarrow$ extract $\rightarrow$ predict $\rightarrow$ generalize) enabling autonomous prediction.

\subsubsection{Summary}

We validated that fix-impact-ratio discriminates strategic from sequential debugging (d=3.07, p=0.0001) and verified the two-stage mechanism: clustering (coefficient 1.909, p=0.001) enables prioritization (42.5\% vs 19.8\%, 2.$15\times$ improvement). However, transfer learning failed (slope ratio 0.82, p=0.504), establishing strategic debugging operates as feedback loop within revealed tests, not pattern-based learning enabling generalization to unseen tests. Three of four hypotheses passed validation (h-e1, h-m1, h-m2 PASS; h-m3 FAIL), confirming partial mechanism while revealing boundary of strategic debugging capability.
